\section{11차시 --- LED와 부저 동기화}
\label{sec:ch11}

본 차시는 빛과 소리 두 출력 채널을 "동시에" 제어하여 더 풍부한 표현을
만든다. 단순 동기 깜빡임에서 시작해, 점점 빨라지는 카운트다운 부저음,
축하 멜로디(도레미파솔), 그리고 4단계 자유 창작 사운드트랙으로 발전한다.
이전 차시에서 따로 다루던 LED와 부저가 이제 하나의 통합된 "무대"로
결합된다.

\subsection{미션 1 --- LED + 부저 동시 제어!}
\label{subsec:ch11-m1}
\missionmeta{LED 5개, 부저}{SYNC}

\begin{storybox}
\sayares{알비, 빛이랑 소리를 동시에 낼 수 있어? 신호등이 깜빡이면서 동시에 소리도 나면 더 확실하잖아!}
\sayalbi{삐릿- 물론이에요! LED를 켤 때 부저도 함께 켜면 빛과 소리가 동시에 나요!}
\end{storybox}

\subsubsection*{학습 목표}
\begin{itemize}[leftmargin=1.4em]
  \item LED 켜기 + 부저 켜기 동시 실행
  \item 0.5초 대기 후 LED 끄기 + 부저 끄기 동시에
  \item 빛과 소리가 완벽하게 동기화된 깜빡임 확인
\end{itemize}

\begin{labbox}{샘플 코드 --- LED + 부저 동기화}
\begin{minted}{python}
# LED + 부저 동기화
for i in range(5):
    led_on(1, 1.0)
    buzzer_on(500, 0.5)
    time.sleep(0.5)
    led_off(1)
\end{minted}
\end{labbox}

\subsection{미션 2 --- 카운트다운 발사 시스템!}
\label{subsec:ch11-m2}
\missionmeta{LED 5개, 부저}{SYNC}

\begin{storybox}
\sayalbi{삐릿! 카운트다운 소리를 만들어보는 건 어떨까요?}
\sayares{좋아! 삐빅 소리가 점점 빨라지면서 출발! 우주 발사 현장을 만들어보자!}
\end{storybox}

\subsubsection*{학습 목표}
\begin{itemize}[leftmargin=1.4em]
  \item LED 하나씩 꺼지며 카운트다운 (5→1)
  \item 각 숫자마다 부저 경고음 함께 출력
  \item 카운트다운이 진행될수록 소리가 점점 빨라짐
  \item 0이 되는 순간 LED 전체 깜빡 + 발사음
\end{itemize}

\begin{labbox}{샘플 코드 --- 카운트다운 발사!}
\begin{minted}{python}
# 카운트다운 발사!
for i in range(1,6): led_on(i,1.0)
freqs = [400,500,600,700,800]
for step in range(5):
    led_off(5-step)
    buzzer_on(freqs[step], 0.3)
    time.sleep(0.7)
# 발사!
buzzer_on(1200, 1.0)
for j in range(5):
    for i in range(1,6): led_on(i,1.0)
    time.sleep(0.15)
    for i in range(1,6): led_off(i)
    time.sleep(0.15)
\end{minted}
\end{labbox}

\subsection{미션 3 --- 발사 성공 세레머니!}
\label{subsec:ch11-m3}
\missionmeta{LED 5개, 부저}{SYNC}

\begin{storybox}
\sayares{발사 성공이야!! 알비! 축하 파티를 열어야 해!}
\sayalbi{삐릿♪♪ 야호!! LED 파티 조명에 부저 축하 멜로디까지! 최고의 세레머니를 만들어봐요!}
\end{storybox}

\subsubsection*{학습 목표}
\begin{itemize}[leftmargin=1.4em]
  \item LED 5개 신나는 파티 패턴 깜빡이기
  \item 부저로 축하 멜로디 연주 (도레미파솔)
  \item LED 파티 조명과 부저 멜로디가 어우러지는 세레머니
\end{itemize}

\begin{labbox}{샘플 코드 --- 발사 성공 세레머니!}
\begin{minted}{python}
# 발사 성공 세레머니!
# 파티 조명
for j in range(8):
    for i in range(1,6): led_on(i,1.0)
    time.sleep(0.15)
    for i in range(1,6): led_off(i)
    time.sleep(0.15)
# 축하 멜로디 도레미파솔
notes = [262,294,330,349,392]
for n in notes:
    buzzer_on(n, 0.3); time.sleep(0.1)
\end{minted}
\end{labbox}

\subsection{미션 4 --- 나만의 발사 사운드트랙!}
\label{subsec:ch11-m4}
\missionmeta{LED 5개, 부저}{SYNC}

\begin{storybox}
\sayares{알비, 지금까지 만든 것들로 완전한 발사 사운드트랙을 만들어보자!}
\sayalbi{삐릿♪ 준비 → 카운트다운 → 발사 → 축하 순서로 LED와 부저를 함께 사용!}
\end{storybox}

\subsubsection*{학습 목표}
\begin{itemize}[leftmargin=1.4em]
  \item 발사 준비: 낮은 음 + LED 순차 점등
  \item 카운트다운: LED 하나씩 꺼짐 + 부저 점점 빨라짐
  \item 발사!: LED 전체 깜빡 + 강한 발사음
  \item 축하: LED 파티 + 부저 멜로디
\end{itemize}

\begin{labbox}{샘플 코드 --- 나만의 발사 사운드트랙}
\begin{minted}{python}
# 나만의 발사 사운드트랙
print("STEP1: 발사 준비")
for i in range(1,6): led_on(i,1.0); buzzer_on(300,0.2); time.sleep(0.3)
print("STEP2: 카운트다운")
for i in range(5,0,-1): led_off(i); buzzer_on(400+i*80,0.3); time.sleep(0.7)
print("STEP3: 발사!!!")
buzzer_on(1200,1.0)
for j in range(6): pass  # LED 전체 깜빡 패턴 자유 작성
print("STEP4: 축하!")
\end{minted}
\end{labbox}

\begin{notebox}{음악 + 시각}
도레미파솔(262, 294, 330, 349, 392Hz)의 주파수는 실제 음악의 C4 옥타브에
해당한다. 학습자는 "코드로 음악을 만든다"는 새로운 표현 방식을
체험하며, 부저가 단순 경보가 아니라 작곡의 도구가 될 수 있음을 발견한다.
\end{notebox}
